\section{Method}
\label{sec:method}

\subsection{EWC formulation and embedded adaptation}
\label{sec:ewc_formulation}

\paragraph{Elastic penalty.} EWC fights catastrophic forgetting by anchoring the parameters deemed
important for what has already been learned \cite{Kirkpatrick2017EWC}. After a task, we keep an
anchor point $\theta^\star$ --- the parameters as they stood at the end of that task --- and an
estimate $F$ of the importance of each coordinate. Training on the next task then minimizes the task
loss augmented with a quadratic pull towards that anchor:
\begin{equation}
\label{eq:ewc_loss}
\mathcal{L}(\theta) \;=\; \mathcal{L}_{\text{task}}(\theta)
\;+\; \frac{\lambda}{2} \sum_{i} F_i \,\bigl(\theta_i - \theta^\star_i\bigr)^2 .
\end{equation}
The coefficient $\lambda$ arbitrates between plasticity and retention. Each coordinate is pulled back
with its own stiffness, $\lambda F_i$: a coordinate of high importance becomes expensive to move, a
coordinate of low importance stays free. The sum runs over \emph{all} parameters of the head, weights
and biases alike. There is a single quadratic term, anchored on the latest $\theta^\star$, rather
than a sum over past tasks: the memory cost of the regularizer is therefore constant, two copies of
the parameters, whatever the number of tasks already seen. On the first task, $F$ and $\theta^\star$
do not exist yet and the penalty is absent.

\paragraph{Estimating importance on PC.} $F$ is the diagonal of the Fisher information matrix,
estimated at the end of a task on that task's data. We use its empirical form, that is, evaluated on
the real labels rather than on labels drawn from the model. The quantity actually accumulated is the
squared gradient of the loss \emph{averaged over a batch}, averaged over the $B$ batches traversed:
\begin{equation}
\label{eq:fisher}
F_i \;=\; \frac{1}{B} \sum_{b=1}^{B}
\left( \frac{\partial \mathcal{L}_b}{\partial \theta_i} \right)^{\!2},
\qquad
\mathcal{L}_b \;=\; \frac{1}{|b|} \sum_{n \in b} \ell\bigl(x_n, y_n; \theta\bigr) .
\end{equation}
The budget is capped at \num{200} samples, which is enough to rank the coordinates against one
another --- only that ranking matters, since $\lambda$ absorbs the global scale. Two writings of the
accumulation coexist in the project. The \emph{online} variant in the sense of
\cite{Schwarz2018Online} maintains a single importance and lets it decay, $F \leftarrow \gamma F +
F^{(t)}$ with $\gamma = \num{0.9}$, which gradually fades the oldest tasks. The head actually ported
to the board instead \emph{replaces} its importance at every consolidation: the anchor always follows
the last consolidated task.

\paragraph{What the board can do, and what it cannot.} The port keeps the penalty of
\eqref{eq:ewc_loss}, whose derivative $\lambda F_i (\theta_i - \theta^\star_i)$ adds to the task
gradient. The embedded update reads, for a weight $W$ of layer $l$:
\begin{equation}
\label{eq:board_sgd}
W_{ji} \;\leftarrow\; W_{ji} \;-\; \eta \Bigl(
\underbrace{\delta_j \, a_i}_{\text{task gradient}}
\;+\; \underbrace{\lambda \, F_{ji} \,\bigl(W_{ji} - W^\star_{ji}\bigr)}_{\text{elastic pull}}
\Bigr),
\qquad \eta = \num{0.01} .
\end{equation}
Consolidation, on the other hand, cannot reproduce \eqref{eq:fisher}: estimating it would require
keeping a batch of labelled samples and running a backward pass over them, neither of which a board
in continuous streaming, with no data memory, has at hand. Importance is therefore approximated by
the \emph{magnitude} of the weight itself, maintained as an exponential moving average:
\begin{equation}
\label{eq:board_fisher}
F_{ji} \;\leftarrow\; \alpha \, F_{ji} \;+\; (1-\alpha) \, W_{ji}^2 ,
\qquad
W^\star_{ji} \;\leftarrow\; W_{ji} ,
\qquad \alpha = \num{0.99} .
\end{equation}
This substitution must be named for what it is: $W^2$ is not the Fisher, it is an importance proxy
that requires no data, no label and no backward pass, and whose cost is one read-write per weight. It
assumes that large-magnitude coordinates are the ones carrying the learned function. Two consequences
follow. The embedded penalty covers the weight matrices only, biases having neither importance nor
anchor stored, which saves a third of the regularization state. And at initialization $F \equiv 0$:
until a consolidation has occurred, the elastic pull is zero and the board learns by plain gradient
descent.

\begin{table}[htbp]
\centering
\caption{Adaptation of the model under embedded constraints. The penalty of \eqref{eq:ewc_loss} is
kept identical; what is simplified is its importance estimator and its optimization regime.}
\label{tab:ewc_port}
\footnotesize
\begin{tabular}{@{}p{30mm}p{35mm}p{50mm}@{}}
\toprule
 & \textbf{PC (PyTorch)} & \textbf{Board (\texttt{ewc\_head.c})} \\
\midrule
Importance estimator & squared gradient, Eq.~\eqref{eq:fisher} & magnitude $W^2$, Eq.~\eqref{eq:board_fisher} \\
Data required        & one labelled batch per task & none \\
Accumulation         & replaced at every consolidation & moving average, $\alpha = \num{0.99}$ \\
Penalty scope        & weights and biases & weights only \\
Optimization         & batches, momentum, several epochs & one sample, no momentum, single pass \\
$\lambda$            & \num{1000} & \num{400} \\
Trigger              & end of task & frame flag, or drift detection \\
Initial state of $F$ & first computed at the end of task $1$ & zero, hence no pull before consolidation \\
\bottomrule
\end{tabular}
\end{table}

One last, purely numerical divergence is worth flagging, as it sheds light on the parity
measurements: the board updates each weight in place and then propagates the error to the previous
layer using the \emph{already modified} weight, where the PC framework uses the pre-step weights
throughout. Under frozen inference, both platforms execute the same function and parity is exact;
under the update regime, this divergence and the unit batch are enough to explain why it no longer
quite is.

\subsection{EWC model and C port}

The model is a multilayer perceptron with an EWC head (\texttt{EWCMlpMulticlass}, $k \to 32 \to 16 \to 2$)
trained incrementally as described in Section~\ref{sec:ewc_formulation}. The head is ported to C (\texttt{ewc\_head.c}) for the
board: both the forward pass and the \gls{SGD} update run on the Cortex-M4 in FP32 (the F439ZI has no NPU).
Weights are exported from the PyTorch checkpoint into a generated C header
(\texttt{model\_weights\_ewc.h}) --- never hand-edited --- so that board and PC share \emph{exactly} the
same parameters.

\subsection{Paired PC$\leftrightarrow$board protocol}

To make the comparison honest, board and PC consume the same sequence: same seed ($\texttt{seed}=42$),
same sample order, same normalization (feature indices and arrays are produced by a single source). Two
protocols are evaluated: \emph{frozen} (inference only, weights fixed) and \emph{online} (inference then
continual update). \emph{Parity} is the prediction-to-prediction agreement rate between the two platforms.

\subsection{Quantization formulation}
\label{sec:quant_formulation}

\paragraph{The mapping and its scale.} Quantizing a tensor means replacing its real values by integers
on a regular grid, plus a scale that maps that grid back to the floating-point world. For a weight on
$b$ bits in signed symmetric representation, with $q_{\max} = 2^{\,b-1} - 1$:
\begin{equation}
\label{eq:quant_map}
Q_{ji} = \mathrm{clip}\!\left( \mathrm{round}\!\left( \frac{W_{ji}}{s_j} \right),\, -q_{\max},\, q_{\max} \right),
\qquad
\widehat{W}_{ji} = s_j \, Q_{ji} .
\end{equation}
The whole behaviour of a scheme lies in the choice of $s_j$. We compare three of them, which form a
single family rather than separate methods:
\begin{equation}
\label{eq:scale_gran}
s_j \;=\;
\underbrace{\frac{1}{128}}_{\text{fixed}}
\qquad\text{or}\qquad
\underbrace{\frac{\max_{j,i} |W_{ji}|}{q_{\max}}}_{\text{per tensor}}
\qquad\text{or}\qquad
\underbrace{\frac{\max_i |W_{ji}|}{q_{\max}}}_{\text{per output channel}} .
\end{equation}
The fixed scale is the one of the original embedded kernel. It does not depend on the weights: any
coefficient with modulus above $127/128$ is clamped onto the edge of the grid, and every coefficient
smaller than the step collapses onto zero.

\paragraph{The integer layer.} Activations receive the same treatment, with a scale $s_a$ obtained from
a calibration. A layer then reads as an integer accumulation followed by a single return to floating
point:
\begin{equation}
\label{eq:int_layer}
\mathrm{acc}_j = \sum_i Q_{ji} \, q^a_i ,
\qquad
y_j = \mathrm{acc}_j \cdot s_j \cdot s_a + b_j .
\end{equation}
Granularity is read off \eqref{eq:int_layer}: the index $j$ carried by $s_j$ means that an output neuron
whose weights are of small amplitude gets its own scale, instead of sharing the one set by the largest
coefficient of the whole matrix. The calibration used is a per-layer absolute maximum, measured on a
batch in floating-point precision ($s_a = \mathrm{act\_max}/q^a_{\max}$).

\paragraph{What the original scheme prescribes, and what we keep of it.} The reference integer
formulation \cite{Jacob2018,Krishnamoorthi2018} aims at \emph{fully} integer inference: the return to
the next scale is done there by an integer multiplier and a shift, biases are carried in \texttt{int32}
at the scale of the product, and a zero point offsets the activation grid. Our target has different
constraints: the Cortex-M4F has a hardware floating-point unit, and the head has only three layers. We
therefore keep the integer accumulation and the per-channel scales, but return to floating point
between layers. Table~\ref{tab:quant_ref} summarizes the divergences.

\begin{table}[htbp]
\centering
\caption{Reference integer scheme and embedded implementation. The divergences are dictated by the
target (floating-point unit present, short network) and not by a convenience simplification.}
\label{tab:quant_ref}
\footnotesize
\begin{tabular}{@{}p{30mm}p{40mm}p{45mm}@{}}
\toprule
 & \textbf{Reference \cite{Jacob2018}} & \textbf{v2 kernel (\texttt{ewc\_head\_int8\_v2.c})} \\
\midrule
Scale return      & integer multiplier and shift & floating-point dequantization on the hardware unit \\
Biases            & \texttt{int32} at the product scale & exact floats, never quantized \\
Zero point        & on activations and weights & symmetric weights with no zero point; symmetric activations, affine optional \\
Granularity       & per tensor, per channel as an extension & per output channel by default \\
Calibration       & histogram or percentile & per-layer absolute maximum on a batch \\
Normalization folding & required before quantization & not applicable, the network has none \\
Accumulator       & \texttt{int32} & \texttt{int32}, widened to \texttt{int64} for the Q15 variant \\
\bottomrule
\end{tabular}
\end{table}

\paragraph{The format axis: the ablation ladder.} The original embedded kernel accumulates four
divergences from \eqref{eq:quant_map}--\eqref{eq:int_layer}: a fixed scale, rounding replaced by
truncation towards zero, an \texttt{int16} accumulator that silently wraps, and activations folded onto
a Q7 grid. Rather than fixing them all at once, we walk a trajectory inside family
\eqref{eq:scale_gran}, each step enabling one change only:
\begin{itemize}
  \item \texttt{legacy\_c}: fixed scale, truncation, \texttt{int16} accumulator (the original v1 kernel
        on the board);
  \item \texttt{fix\_acc32}: accumulator widened to \texttt{int32}, everything else unchanged;
  \item \texttt{per\_tensor\_calib}: \emph{calibrated} per-tensor scale and round-to-nearest;
  \item \texttt{per\_channel\_int8}: \emph{per-output-channel} scale;
  \item \texttt{q15}: 16-bit grid, $q_{\max} = 32767$ (RAM $\div 2$ instead of $\div 4$).
\end{itemize}
The v2 firmware kernel implements the last two steps, with an optional Q15 variant
(\texttt{-DEWC\_INT8\_Q15}); the v1 path is unchanged so that both can be compared on the same board.
The ablation is \emph{bit-exactly emulated on PC} (\texttt{src/utils/int8\_c\_emulation.py}) so as to
isolate the contribution of each fix independently of board access. Scales are not recomputed on the
board: they are produced on PC by the very functions the emulator uses, then written into a generated C
header, which makes parity true \emph{by construction} rather than observed after the fact.

\paragraph{Where quantization stops under the learning regime.} A scope statement must be made before
reading the "online" cells of Section~\ref{sec:results}, because it governs their interpretation. The
integer kernel exposes a forward pass only: there is no gradient descent in the quantized domain. The
board therefore keeps a \emph{master} FP32 head, to which \eqref{eq:board_sgd} and
\eqref{eq:board_fisher} apply, and whose integer representation is only a \emph{view}, rebuilt through
\eqref{eq:quant_map} after each step. Two consequences follow, and we state them rather than leave them
implicit. First, the fourfold reduction of weight memory holds for inference alone: under the learning
regime the board simultaneously carries the floating head, its integer view, the importance and the
anchor, that is more than FP32 alone. Second, requantizing every matrix at each sample enters the latency
budget, which Section~\ref{sec:recovery-board} quantifies. In other words, this work establishes a
quantization \emph{compatible} with incremental learning --- inference is integer, and learning keeps
working through it --- not learning \emph{carried out} in integer arithmetic, which remains an open
problem.

\paragraph{What per-channel costs.} A scale per output neuron is not free: it adds an array of floats
parallel to each matrix, one four-byte value per output row. On the ported head, these scales bring the
actual weight-memory gain below the factor of four advertised by the type change alone. That is the
price of robustness: per-channel isolates the large-amplitude neurons, which would otherwise impose
their scale on the whole layer.

\subsection{The moment axis: when to quantize}

At a fixed format, the same head can be quantized at three instants. \emph{Before} training, through QAT:
a \emph{fake-quant} inserted in the loop makes the weights learn under the reduced-precision constraint.
\emph{After}, through PTQ: the trained FP32 model is converted once and for all. And \emph{both}, that is
a QAT training whose weights are then exported to the calibrated integer kernel. Only this last path is
faithful to the actual deployment, since it is the embedded kernel that runs inference, not the
quantization simulation of the training framework. The three moments are compared at identical model,
dataset and seed.

\subsection{The depth axis: bits, granularity, symmetry}

Below INT8, three parameters describe a scheme: the number of \emph{weight bits} (4, 2, ternary, binary),
the \emph{granularity} of the scale (per tensor or per output channel) and the \emph{symmetry}
(symmetric, or affine with a zero point). The first two are read directly off \eqref{eq:quant_map} and
\eqref{eq:scale_gran}: going deeper means lowering $q_{\max}$, and changing granularity means changing
the scope of the maximum that defines $s_j$.

At very low depths, the linear grid gives way to two dedicated schemes, which we take as they stand. In
ternary \cite{Li2016TWN}, a per-channel threshold decides which weights survive, and the scale is the
mean of the surviving weights alone:
\begin{equation}
\label{eq:twn}
\begin{gathered}
\Delta_j = \num{0.7} \cdot \overline{|W_{j\cdot}|},
\qquad
Q_{ji} = \mathrm{sign}(W_{ji}) \cdot \mathbf{1}\bigl[\,|W_{ji}| > \Delta_j\,\bigr], \\[2pt]
\alpha_j = \overline{\bigl\{ |W_{ji}| : |W_{ji}| > \Delta_j \bigr\}} .
\end{gathered}
\end{equation}
In binary \cite{Rastegari2016XNOR}, no weight is set to zero any more and the scale is the mean modulus
of the channel: $Q_{ji} = \mathrm{sign}(W_{ji})$, $\alpha_j = \overline{|W_{j\cdot}|}$. Both schemes
carry a per-channel scale \emph{by construction}, so the granularity axis stops applying once they are
reached.

Symmetry, for its part, concerns activations. After a positive activation function, the negative half of
a symmetric grid is of no use, which costs all the more as bits are scarce. The affine variant offsets
the grid by a zero point:
\begin{equation}
\label{eq:affine_act}
q = \mathrm{round}\!\left(\frac{a}{s_a}\right) + z,
\qquad
\widehat{a} = (q - z)\, s_a .
\end{equation}
Weights, whose distribution is centred, stay symmetric with no zero point in every scheme we evaluate.

One point finally deserves to be stated, as it governs how the
advertised gains must be read: as long as weights remain stored in an \texttt{int8\_t} container, reducing
the number of bits frees \emph{no} byte. The memory gain exists only with effective \emph{bit-packing},
which in turn costs unpacking at run time. We therefore systematically distinguish the \emph{theoretical}
RAM of a scheme from the \emph{measured} RAM of the binary.

\subsection{Measurement methodology}

The quantities reported in Section~\ref{sec:system} call for three distinct instrumentations.

\paragraph{Memory.} The RAM of a binary is not the footprint of its weights. We report \emph{total RAM}
$= \texttt{.data} + \texttt{.bss} + \text{stack peak}$, the first two terms being read from the binary and
the third measured at run time by stack painting, keeping the maximum between the inference phase and the
update phase.

\paragraph{Latency.} The \gls{DWT} cycle counter times the kernel to the cycle. For the INT8 path it is
instrumented in \emph{three segments} --- dequantization, integer \gls{MAC}, requantization --- so as to answer
not "is INT8 slower?" but "which item costs?". Cycles are reported raw, since conversion to microseconds
truncates the information at this scale.

\paragraph{Energy.} Current is measured with an X-NUCLEO-LPM01A probe through three independent methods:
regression of current against inference rate, difference against a \gls{WFI} idle state, and a batch of
inferences per frame isolating computation from the link cost. We report them \emph{side by side} without
ever averaging them: they do not measure the same boundary, and their disagreement is itself a result. Any
unmeasured cell is reported as "to be measured" rather than estimated.
